\section{Scalar transitions and physical section calculus}

The regional passages require estimates on normalized scalar maps, on
derivatives with moving initial sections, and on local Dulac maps. We
develop those tools here. Regular adapters and divided closing maps are
handled by the identities of Section~\ref{sec:counting-tools}.

\subsection{Normalized scalar transitions}
Subtracting a reference orbit and dividing by its first variation removes
the affine part of a scalar transition. The remaining equation vanishes
to second order in the normalized phase variable. An integrable bound on that
equation gives uniform endpoint maps even when the unnormalized first
variation has no uniform bound.

\begin{lemma}[Normalized scalar transitions]\label{lem:scalar-normalization}
Let \(k\ge1\). Consider equations
\[
 \eta_r=N_\vartheta(r,\eta),\qquad L\le r\le R_\vartheta,
 \qquad N_\vartheta(r,0)=\partial_\eta N_\vartheta(r,0)=0,
\]
on \(|\eta|\le a\), where \(L\) and \(a>0\) are fixed.
Suppose that, uniformly in the parameters,
\[
 \|N_\vartheta(r,\cdot)\|_{C^{k+1}([-a,a])}\le g(r),
 \qquad g\in L^1([L,\infty)).
\]
Then there is a fixed neighborhood of zero on which all endpoint
maps have bounded \(C^{k+1}\) norms and positive first derivatives
bounded away from zero. Their inverses have the same bounds on a
fixed smaller neighborhood of zero.

Suppose also that \(R_\vartheta\to\infty\) and the kernels converge
in \(C^0([L,M];C^{k+1}_\eta)\) for every fixed \(M>L\).
Then the endpoint maps
and their inverses converge in \(C^{k+1}\) on smaller neighborhoods.
For a chosen parameter derivative \(\partial_\ell\), convergence
in \(C^k\) holds for the differentiated maps if the differentiated
kernels converge in \(C^0([L,M];C^k_\eta)\), have a common integrable
majorant, and
\[
 \|N_\vartheta(R_\vartheta,\cdot)
               \partial_\ell R_\vartheta\|_{C^k}\longrightarrow0.
\]
The convergence is joint in an additional compact parameter family
if these local convergences, \(R_\vartheta\to\infty\), and the
vanishing endpoint term are uniform in that family.
\end{lemma}
\begin{proof}
Taylor's formula gives \(|N_\vartheta(r,\eta)|\le g(r)|\eta|^2/2\).
Comparison with \(y'=g(r)y^2/2\) keeps all solutions from a common
enlarged initial interval inside \(|\eta|<a\). The first variation is
\[
 \exp\left(\int_L^r
          \partial_\eta N_\vartheta(t,\eta(t))\,dt\right)>0.
\]
It and its reciprocal are bounded. Induction in the triangular
variational equations and Gronwall's inequality bound all remaining
rows through order \(k+1\). The slope lower bound gives a common
image interval, and the inverse identity gives the inverse bounds.

Subtracting the integral equations and their variations gives convergence
on \([L,M]\); each remaining tail costs at most \(C\int_M^\infty g\).
Let the parameters tend to their limit and then \(M\to\infty\).
For the parameter derivative, the variation satisfies
\[
 Z_r=\partial_\eta N_\vartheta\,Z+\partial_\ell N_\vartheta,
 \qquad Z(L)=0.
\]
The endpoint derivative adds
\(N_\vartheta(R_\vartheta,\eta(R_\vartheta))
\partial_\ell R_\vartheta\). The same argument proves \(C^k\)
convergence under the endpoint assumption: differentiating a coefficient
\(N_{\eta^h}(r,\eta)\) adds
\(N_{\eta^{h+1}}(r,\eta)Z\), which requires the extra phase order.
Finally, \(C^{k+1}\) convergence of the endpoint maps \(A_\vartheta\) and
\[
 \partial_\ell A_\vartheta^{-1}
 =-\left(\frac{\partial_\ell A_\vartheta}{A_\vartheta'}\right)
     \circ A_\vartheta^{-1}
\]
give \(C^k\) convergence of the differentiated inverses. Each estimate
is uniform in the additional compact family under the stated assumptions.
\end{proof}

For a graph equation \(Y_r=F(r,Y)\), a reference solution \(Y_*\)
and its first variation \(j>0\), the substitution
\(\eta=(Y-Y_*)/j\) gives the exact kernel
\begin{equation}\label{tail:normalized-kernel}
 N(r,\eta)=j\eta^2\int_0^1(1-t)
                  F_{YY}(r,Y_*+tj\eta)\,dt.
\end{equation}
For \(h\ge2\), \(\partial_\eta^hN=j^{h-1}F_{Y^h}(r,Y_*+j\eta)\).
At fixed \(r,\eta\), a parameter derivative is therefore
\begin{equation}\label{deg:det-parameter-phase-derivative}
\begin{split}
 \partial_\ell\partial_\eta^hN
 =j^{h-1}\bigl\{&(h-1)(\partial_\ell\log j)F_{Y^h}
 +\partial_\ell^{\rm exp}F_{Y^h}\\
 &+F_{Y^{h+1}}(Y_{*,\ell}+j_\ell\eta)\bigr\}.
\end{split}
\end{equation}
Here \(\mathrm{exp}\) holds the graph coordinates fixed. The orders
zero and one follow by integrating in \(\eta\), since
\(N(r,0)=N_\eta(r,0)=0\). Thus \(C^3\)-valued parameter
derivatives require \(F_{Y^4}\) and \(N_{\eta^4}\).
Integrating the indicated parameter derivative gives the sequential
parameter secants used below. This requires neither a mixed parameter
derivative nor convergence of \(j\).

A quantitative form controls how close the normalized map is to the
identity when only an upper bound on the first variation is available.
For an actual graph family
write \(J_n=\partial_b^nY\), \(Q_n=J_n/J_1\), and
\(a_n=F_{Y^n}J_1^{n-1}\). Direct differentiation gives
\begin{equation}\label{tail:relative-jet-equations}
 Q_2'=a_2,\qquad Q_3'=3a_2Q_2+a_3,\qquad
 Q_4'=a_2(4Q_3+3Q_2^2)+6a_3Q_2+a_4.
\end{equation}
If, uniformly on an enlarged label interval,
\(Q_n(L)=O(\varepsilon)\) and
\(\int|a_n|\,dr=O(\varepsilon)\), \(2\le n\le4\),
these identities give \(Q_n=O(\varepsilon)\).
Since \(\partial_b\log J_1=Q_2\), the endpoint map in \(b-b_*\),
normalized by its exact value and slope at \(b_*\), and its inverse on a
smaller common interval, are \(\mathrm{id}+O_{C^3}(\varepsilon)\)
for sufficiently small \(\varepsilon\).
This conclusion is independent of the size of \(J_1(b_*)>0\).


\subsection{Moving-section variations and inward continuation}

Moving a section changes both the initial height and the time of
departure. The following identities retain both effects, including the
endpoint terms that appear after using height as the independent variable.

\begin{lemma}[Scalar differentiation on moving sections]
\label{deg:lem-scalar-moving-identities}
Fix the coefficients of a scalar graph \(z_x=f(x,z)\), with initial
curve \((x_l(a),z_l(a))\), and let \(D\) be a phase vector field,
commuting with differentiation at fixed \(x\). Write \(V_j=D^jz\).
Through order three its equations and moving initial rows are
\begin{equation}\label{deg:scalar-moving-bell}
 \begin{gathered}
 (V_1)_x=f_zV_1,\qquad (V_2)_x=f_zV_2+f_{zz}V_1^2,\\
 (V_3)_x=f_zV_3+3f_{zz}V_1V_2+f_{zzz}V_1^3,\\
 B_1=Dz_l-f_lDx_l,\quad B_2=DB_1-f_{z,l}B_1Dx_l,\\
 B_3=DB_2-(f_{z,l}B_2+f_{zz,l}B_1^2)Dx_l,
 \qquad V_j(x_l;a)=B_j .
 \end{gathered}
\end{equation}
For \(f=z(n+\zeta z)/(M+Qz)\) on a positive graph with positive
denominator, put \(\mathscr D=M+Qz\).
The complete fixed-port first derivative, including an observer \(\chi\), is
\begin{equation}\label{deg:scalar-raw-cocycle}
 \begin{gathered}
 f_z=\frac d{dx}\log\frac z{\mathscr D}
          +\frac{M'+(Q'+\zeta)z}{\mathscr D},\\
 D\chi(z_T)=\chi'(z_T)B_1
    \frac{(z/\mathscr D)_T}{(z/\mathscr D)_l}
    \exp\int_{x_l}^{x_T}\frac{M'+(Q'+\zeta)z}{\mathscr D}\,dx .
 \end{gathered}
\end{equation}
For \(z_l=\ell_0M(a)\), \(n=xB\), and
\(b=-a^2/(2\sigma^2)\), this becomes
\begin{equation}\label{deg:scalar-polynomial-lip-slope}
 \begin{gathered}
 \Gamma_b'=\sigma^2\chi'(z_T)\frac{z_T}{\mathscr D_T}B(a)K(a)e^{I(a)},
 \quad I=\int_a^{x_T}\frac{M'+(Q'+\zeta)z}{\mathscr D}\,dx,\\
 K=1-\frac{M'(a)}{aB(a)}(1+\ell_0Q(a))
       +\frac{\ell_0\zeta(a)M(a)}{aB(a)}.
 \end{gathered}
\end{equation}
If instead \(r_\eta=V(r,\eta)\), with moving data
\((\eta_l,r_l)\), write \(P_j=D^jr|_\eta\), \(l_j=D^j\eta_l\),
\(t_j=D^j\eta_T\), where \(r_T\) is fixed. Then
\begin{equation}\label{deg:scalar-inverse-endpoints}
 \begin{gathered}
 P_{1,l}=Dr_l-V_ll_1,\quad
 P_{2,l}=DP_{1,l}-V_{r,l}P_{1,l}l_1,\\
 (P_1)_\eta=V_rP_1,\quad
 (P_2)_\eta=V_rP_2+V_{rr}P_1^2,\\
 t_1=-P_{1,T}/V_T,\quad
 t_2=-\{P_{2,T}+2V_{r,T}P_{1,T}t_1
              +(V_{r,T}V_T+V_{\eta,T})t_1^2\}/V_T .
 \end{gathered}
\end{equation}
For \(I=\int_{\eta_l}^{\eta_T}J(r(a,\eta),\eta)\,d\eta\),
its full rows are
\begin{equation}\label{deg:scalar-density-endpoints}
 \begin{aligned}
 DI={}&J_Tt_1-J_ll_1+\int J_rP_1\,d\eta,\\
 D^2I={}&\int(J_{rr}P_1^2+J_rP_2)\,d\eta
       +(J_rP_1)_Tt_1-(J_rP_1)_ll_1\\
 &+J_Tt_2+J_{\eta,T}t_1^2-J_ll_2
       -(J_{r,l}Dr_l+J_{\eta,l}l_1)l_1 .
 \end{aligned}
\end{equation}
The observer and divided-closing rows use
Lemma~\ref{phys:lem-c3-adapters}, specifically
\eqref{phys:eq-observer-log} and \eqref{phys:eq-closing-jets}.
Coefficient-only gains and affine calibrations are frozen under \(D\).
\end{lemma}
\begin{proof}
Differentiate the graph equation at fixed \(x\); differentiating
\(V_j(x_l(a),a)=B_j(a)\) gives the initial recursion.
Rational differentiation proves \eqref{deg:scalar-raw-cocycle}.
For the inverse graph differentiate twice at fixed \(\eta\), and
then twice the identity \(r(a,\eta_T(a))=r_T\).
Leibniz differentiation gives \eqref{deg:scalar-density-endpoints};
at the lower endpoint the total radial derivative is \(Dr_l\),
whereas at the upper endpoint it is zero. The observer and closing
claims follow from the cited ordinary \(C^3\) identities.
\end{proof}

In the inward passages, the ratio of height to the axis polynomial
decays on the action clock. The next lemma estimates the resulting
integral without discarding the dependence on its terminal clock time.

\begin{lemma}[Stable ratio integrals with moving endpoints]
\label{deg:lem-stable-ratio-ledger}
At fixed action clock \(q\), let a phase field \(D\) commute with
\(\partial_q\). For \(\kappa>0\), consider
\[
 W_q=\mathcal F=W\{- (1+\omega W)/(1+QW)+\kappa H\},
 \qquad W(0)=\ell_0,
 \qquad L=(Q-\omega)W/(1+QW).
\]
Suppose \(Q,\kappa H,\omega\) and their required two clock rows are
uniformly bounded, \(1+QW\) is a positive unit, and both
\(\mathcal F/W\) and \(\mathcal F_W\) are at most \(-c_0<0\)
on \(0\le W\le\ell_0\). For some \(0\le\rho\le\rho_0\),
assume the first three phase rows of \(\kappa H,Q,\omega\),
and mixed rows containing a phase derivative of total order at most three,
are bounded by
\(C\rho e^{\alpha q}\), with \(\alpha<c_0/12\).
There are an integer \(N\) and \(0<c<c_0/2\), depending only on the
stated bounds, such that
\begin{equation}\label{deg:stable-ratio-fixed-clock}
 0<W\le\ell_0e^{-c_0q},\qquad
 |D^jW|+|D^jL|\le C\rho(1+q)^Ne^{-cq},\quad1\le j\le3.
\end{equation}
The needed mixed rows obey these bounds; pure clock rows have the
same exponential bound without \(\rho\).
For a moving \(T\), put \(T_j=D^jT\) and \(I_T=\int_0^T L\,dq\).
All endpoint contributions are
\begin{equation}\label{deg:stable-ratio-terminal-rows}
 \begin{aligned}
 DI_T={}&\int_0^T DL\,dq+L_TT_1,\\
 D^2I_T={}&\int_0^T D^2L\,dq+2(DL)_TT_1+(L_q)_TT_1^2+L_TT_2,\\
 D^3I_T={}&\int_0^T D^3L\,dq+3(D^2L)_TT_1
       +3(DL_q)_TT_1^2+(L_{qq})_TT_1^3\\
 &+3(DL)_TT_2+3(L_q)_TT_1T_2+L_TT_3 .
 \end{aligned}
\end{equation}
If \(v_q=-\kappa P(v)/(vB(v))\), \(H=P'/(vB)\), and
\(W=Y/P(v)\), its exact polynomial memory is
\begin{equation}\label{deg:stable-ratio-polynomial-memory}
 W(q)=\ell_0\frac{P(a)}{P(v(q))}e^{-q+I(q)},\qquad
 Y_T=\ell_0P(a)e^{-T+I_T}=P(v_T)W(T).
\end{equation}
The corresponding rows of \(W(T)\) are the ordinary chain rule
\(DW=W_D+W_qT_1\),
\(D^2W=W_{DD}+2W_{Dq}T_1+W_{qq}T_1^2+W_qT_2\), and its
third row with coefficients \(1,3,3,1,3,3,1\).
Thus the interior integral contributions are \(O(\rho)\); the full
moving-limit rows also retain their terminal contributions, which are at most
\(C(1+T)^Ne^{-cT}(1+\sum_{j=1}^3|T_j|)^3\).
The constants depend only on the stated coefficient bounds,
\(c_0,\ell_0,\rho_0\), independently of the flight length.
\end{lemma}
\begin{proof}
The negative logarithmic rate gives the first-exit bound for \(W\).
Differentiate its equation three times. Each row has diagonal
\(\mathcal F_W\); the explicit coefficient forcing is
\(O(\rho e^{\alpha q}W)\), and other terms are products of lower
rows. Successive variation of constants absorbs at most three
coefficient growth factors into the fixed exponential decay and
proves \eqref{deg:stable-ratio-fixed-clock}. Differentiating these
equations in \(q\) gives the mixed rows. Leibniz differentiation and
the ordinary chain rule give the terminal formulas.
Since \(\partial_q\log P(v)=-\kappa H\), integrating the ratio
equation gives \eqref{deg:stable-ratio-polynomial-memory}.
\end{proof}

When the radial route is fixed, the same decay controls ordinary
derivatives with respect to the initial height, even for arbitrarily
long flights.

\begin{lemma}[Initial-height rows on a frozen ratio route]
\label{deg:lem-fixed-route-height-rows}
Consider the ratio equation and density \(\mathcal F,L\) of
Lemma~\ref{deg:lem-stable-ratio-ledger}, with coefficients and radial
endpoints fixed. On its full prescribed buffer \(0\le W\le\ell_0\),
assume the stated positive-unit denominator, the two negative margins
\(\mathcal F/W,\mathcal F_W\le-c_0<0\), and bounded
\(Q,\omega,\kappa H\). Let \(W(0;w)=w\), and, for a fixed finite
flight \(T\ge0\), put
\[
 \mathcal N_T(w)=\int_0^T L(q,W(q;w))\,dq,\qquad e_w=w\partial_w.
\]
Uniformly in \(T\) and \(0\le w\le\ell_0\), the ordinary extension
at \(w=0\) satisfies
\[
 \mathcal N_T(0)=0,\qquad
 |e_w^j\mathcal N_T(w)|\le Cw\quad(0\le j\le3),\qquad
 \sum_{j=1}^3|\partial_w^j\mathcal N_T(w)|\le C.
\]
The normalized nonlinear unit \(e^{\mathcal N_T}\) and its reciprocal
have bounded ordinary \(C^3\) norms. In the polynomial-memory setting,
\(Y_T=Y_0e^{-T+\mathcal N_T(w)}\), with \(Y_0=P(v_0)w\).
\end{lemma}
\begin{proof}
The buffer and negative logarithmic margin give
\(0\le W(q;w)\le we^{-c_0q}\). Rational height differentiation
bounds \(\mathcal F_{WW},\mathcal F_{WWW}\) and the first three
height derivatives of \(L\). Apply the triangular equations
\eqref{deg:scalar-moving-bell} to \(V_j=\partial_w^jW\), with
\(V_1(0)=1\), \(V_2(0)=V_3(0)=0\). Successive variation of
constants gives \(|V_j(q)|\le Ce^{-c_0q}\), \(1\le j\le3\).
Every term in \(\partial_w^jL(q,W(q;w))\) contains a decaying
variation, so its integral is uniformly bounded. Also \(|L|\le CW\)
gives \(|\mathcal N_T|\le Cw\). The ordinary extension at zero and
the expansion of \((w\partial_w)^j\) give the weighted rows.
Exponentiation and
\eqref{deg:stable-ratio-polynomial-memory} prove the final claims.
\end{proof}

After the long inward flight, a short central segment remains. Factoring
out its linear propagator gives a nonlinear error controlled by the
incoming height, uniformly as that propagator tends to zero.

\begin{lemma}[Compact inward height continuation]
\label{deg:lem-compact-height-continuation}
For a frozen coefficient field, suppose the actual equation on a
fixed segment, oriented from \(b>0\) to zero, is
\(Y_v=A(v)Y+K(v,Y)Y^2\), with
linear propagator \(m(b)=1\), \(0<m(v)\le1\). On its prescribed
height buffer \(0\le Y\le h_0\), assume the ordinary height
derivatives of \(K\) through order three are bounded by
\(C/\kappa_c\), for a local parameter \(\kappa_c>0\).
For \(0<y\le h_0/2\) and \(y/\kappa_c\) sufficiently small,
its actual continuation satisfies
\begin{equation}\label{deg:scalar-compact-propagator}
 Y(0;y)=m(0)y e^{\mathcal N(y)},\qquad
 |(y\partial_y)^j\mathcal N(y)|\le Cy/\kappa_c,\quad0\le j\le3.
\end{equation}
In composition with \(y=y(a)\), with the field, segment and
\(\kappa_c\) fixed, put \(\Lambda_j=D^j\log y\) and
\(e=y\partial_y\). The three rows of \(\mathcal N(y(a))\) are
\(e\mathcal N\,\Lambda_1\),
\(e^2\mathcal N\,\Lambda_1^2+e\mathcal N\,\Lambda_2\), and
\(e^3\mathcal N\,\Lambda_1^3+3e^2\mathcal N\,\Lambda_1\Lambda_2
+e\mathcal N\,\Lambda_3\). No reciprocal bound on \(m(0)\) is used.
Constants depend only on the stated height bounds and fixed segment,
independently of the extracted multiplier.
\end{lemma}
\begin{proof}
Write \(Y=m(v)yu(v;y)\). The integral
equation for \(u\), with \(u(b)=1\), has coefficients and its first
three Euler forcings bounded by \(Cy/\kappa_c\) on
\(1/2\le u\le2\). First exit and three triangular integral
estimates give \(u^{\pm1}\le C\) and
\(|e^j\log u(0;y)|\le Cy/\kappa_c\). The displayed composition
rows follow by the chain rule.
\end{proof}

\subsection{An exterior spectral estimate}

The outward cubic, quadratic and degree-loss passages share a scalar
graph equation. The quadratic coefficient determines a power in its
first variation, while the remaining terms control the error in its
logarithmic derivatives. We use logarithmic height as the independent
variable to estimate these derivatives directly, including the moving
endpoints.

\begin{lemma}[Exterior spectral estimate]\label{deg:lem-spectral-exterior}
Fix a terminal cut \(d>0\), a small lip constant \(\ell_0>0\), compact
bounds on \(A,C\), and positive constants \(c_0,C_0\). Put
\[
 \begin{gathered}
 M=tx^3+\mu_2x^2+\mu_1x+\mu_0,\qquad B=1+b+tx,\\
 Q=1-A(v+tx)+C(v+tx)^2,\quad \zeta=tC(v+tx),\quad |t|\le1.
 \end{gathered}
\]
All coefficients are fixed under \(E=a\partial_a\). On an enlarged
open interval of \(0<a<d/2\), suppose, pointwise in \(a\), that
some \(0<\rho\le\rho_0\) satisfies
\begin{equation}\label{deg:spectral-exterior-hypotheses}
 \begin{gathered}
 c_0\rho\le m(x):=M(x)/x^2\le C_0\rho\quad(a\le x\le2a),\\
 |b|+|v|+|\mu_2|+|t|a+|\mu_1|/a+|\mu_0|/a^2\le C_0\rho .
 \end{gathered}
\end{equation}
For sufficiently small \(|t|d,\rho_0\), depending only on these
fixed bounds, the graph
\begin{equation}\label{deg:spectral-exterior-graph}
 z_x=\frac{z(xB+\zeta z)}{M+Qz},\qquad z(a)=\ell_0M(a)
\end{equation}
reaches \(x=d\) with positive height and denominator. An ordinary
\(C^3\) observer \(\chi\), with \(\chi'>0\), is appended on the
open preimage of its prescribed buffered terminal image. Set
\[
 \nu=\frac{2\mu_2}{1+b},\quad\lambda=2\nu,\quad e=2-\lambda>3/2,
 \quad\Theta=|t|a+|\mu_1|/a+|\mu_0|/a^2,\quad L_\rho=1+|\log\rho|,
\]
\[
 H_\chi=\sup\{|z(\log\chi')'|
       +|z(\log\chi')'+z^2(\log\chi')''|\},\qquad
 \mathcal T=|t|d+|\mu_2|+|\mu_1|/d+|\mu_0|/d^2+H_\chi ,
\]
where the supremum is on that buffer and is finite.
For any coefficient-only \(\sigma>0\), use
\(\beta=-a^2/(2\sigma^2)\) and \(\Gamma(\beta)=\chi(z_d)\).
Then \(\Gamma'>0\), and
\begin{equation}\label{deg:spectral-exterior-rows}
 |E\log\Gamma'+\lambda|+|E^2\log\Gamma'|
 \le C\{\Theta L_\rho+\mathcal T(a/d)^e\}.
\end{equation}
The same rows hold after a positive coefficient-only affine calibration.
For intrinsic terminal height take \(\chi=\mathrm{id}\), so \(H_\chi=0\).
No reciprocal bound on \(\sigma\) or bound on \(\lambda\log a\) is used.
\end{lemma}
\begin{proof}
Set \(W=z/x^2\), \(\tau=\log(x/a)\), \(N=B+\zeta xW\), and
\(D=m+QW\). Then \(W_\tau=F=W(N/D-2)\). On \([a,2a]\),
\(N,Q\) are positive units and \(D\asymp\rho+W\) for bounded \(W\).
Below a small fixed \(w_0\), \(F\ge cW/(\rho+W)\). Starting at
\(W_l=\ell_0m(a)\), the time to reach \(w_0\) is at most
\[
 C\rho\log\frac{w_0}{\ell_0c_0\rho}+Cw_0<\log2,
\]
by choosing \(w_0\) first and then \(\rho_0\). At the boundaries
\(w_0<1/2\), \(w_1=3/4\), the formula \(F(x,w)=1-2w+O_w(\rho)\)
makes the strip invariant up to \(2a\). Beyond it,
\(|m(x)|\le|t|d+C\rho\ll w_0^2\) preserves the strip, gives
positive-unit \(D\) and \(F_W\le-c<0\). Thus the graph reaches
\(d\), with \(c\rho\le W\le w_1\), even if \(m\) changes sign.

The bounded strong graph with \(b=v=\mu_j=0\), keeping \(t,A,C\),
is \(W_*=1/2+O(|t|x)\), by the integral from zero for
\(xU'=-2U+O(|t|x)\), \(U=W-1/2\). Strip contraction gives
\begin{equation}\label{deg:spectral-exterior-window}
 |W(d)-W_*(d)|\le C\rho+C(a/d)^c.
\end{equation}
Where this comparison applies, the original buffer margin and the
permitted cutoff choices put the terminal height in the observer's
domain. Elsewhere we retain the original terminal admission condition.

Put \(r=\log x\), \(\eta=\log z\). At fixed height, \(W_r=-2W\)
and \(r_\eta=V=D/N\). The raw cocycle
\eqref{deg:scalar-raw-cocycle} has the exact density split
\begin{equation}\label{deg:spectral-exterior-density}
 \begin{gathered}
 \frac{2m+xm'+(xQ'+\zeta x)W}{N}=\nu+J,\\
 J=\frac{(3-\nu)tx+\mu_1/x+
          tx[-A+(3-\nu)C(v+tx)]W}{N}.
 \end{gathered}
\end{equation}
Here \(2m+xm'=3tx+2\mu_2+\mu_1/x\); \(\mu_0\) cancels only
in the density. Let
\(I=\int_{\eta_l}^{\eta_T}J\,d\eta\), with
\(\eta_l=\log(\ell_0M(a))\). The near height length is \(O(L_\rho)\);
on the strip \(d\eta\le C\,d\log x\). Thus every partial integral
is bounded by \(C\{\Theta L_\rho+|t|d+|\mu_1|/a\}\).

Use \eqref{deg:scalar-inverse-endpoints} with
\(P_j=E^jr|_\eta\), \(l_j=E^j\eta_l\), \(t_j=E^j\eta_T\),
\(Er_l=1\). The polynomial gives bounded \(l_1,l_2\),
\(V_l,V_{r,l}=O(\rho)\), and \(P_{1,l}=1+O(\rho)>0\).
For \(\mathscr B=EW_l-F_l\), the exact first cocycle gives
\begin{equation}\label{deg:spectral-exterior-inverse-decay}
 \begin{gathered}
 P_1=\frac{-\mathscr B(1+\ell_0Q(a))}{\ell_0N}
       (a/x)^e\left(\frac W{\ell_0m(a)}\right)^\nu e^I,\\
 0<P_1\le C(a/x)^e,\qquad |P_2|\le CP_1 .
 \end{gathered}
\end{equation}
The initial factor equals \(B(a)K(a)>0\), with \(K\) from
\eqref{deg:scalar-polynomial-lip-slope}. Since \(|\nu|\le C\rho\)
and \(c\rho\le W\le w_1\), its height ratio to power \(\nu\) is a
unit for either sign, proving the first bound without bounding
\(a^{-\lambda}\). Near the lip, fixed-height differentiation gives
\(|V_{rr}|\le C(\rho+W)\), whereas \(d\eta=N/D\,d\log x\).
On the strip all these factors are bounded. Hence
\[
 |V_{rr}|\,d\eta\le C\,d\log x,\qquad
 (P_2/P_1)_\eta=V_{rr}P_1.
\]
Its integral is bounded by \(C\int_a^d(a/x)^e\,d\log x\), proving
the second bound. The two fixed-terminal rows then give
\(|t_1|+|t_2|\le C(a/d)^e\).

Rational differentiation at fixed height gives
\(|J_r|+|J_{rr}|\le C(|t|x+|\mu_1|/x)\). In the two full Leibniz
rows \eqref{deg:scalar-density-endpoints}, the lower terms cost
\(C\Theta\), the near integrals cost \(C\Theta L_\rho\), and
\[
 \int_{2a}^d(|t|x+|\mu_1|/x)(a/x)^e\,d\log x
 \le C(|t|a+|\mu_1|/a).
\]
Every upper term retains \((a/d)^e\). Therefore
\begin{equation}\label{deg:spectral-exterior-integral-rows}
 |EI|+|E^2I|
 \le C\{\Theta L_\rho+(|t|d+|\mu_1|/d)(a/d)^e\}.
\end{equation}
The exact polynomial-lip formula becomes
\begin{equation}\label{deg:spectral-exterior-slope}
 \Gamma'=\sigma^2a^{-\lambda}
 B(a)K(a)\ell_0^{-\nu}m(a)^{-\nu}
 \frac{W_T}{m(d)+Q(d)W_T}z_T^\nu\chi'(z_T)e^I .
\end{equation}
Both \(E\) rows of \(\log(BK)\) are \(O(\Theta)\). Also
\(E m,E^2m=O(\Theta)\), \(m(a)\asymp\rho\), and
\(|\nu|,\Theta\le C\rho\), so those of \(-\nu\log m(a)\) cost
\(C\Theta\). At fixed \(d\), the first two height-logarithmic rows
of \(\log[W_T/(m(d)+Q(d)W_T)]\) cost \(C|m(d)|\), and those of
\(\nu\log z_T\) cost \(C|\nu|\). The observer rows are precisely
those defining \(H_\chi\). Their terminal chain rules and the integral
bounds prove \eqref{deg:spectral-exterior-rows}, since
\(a^{-\lambda}\) contributes \(-\lambda\) and zero to the first and
second Euler rows.
\end{proof}


\subsection{Exact identities on moving polynomial sections}
\label{subsec:deg-moving-sections}

For the cubic, quadratic and degree-drop scalings consider
\begin{equation}
 \dot\xi=-\epsilon\{P(\xi)+Q(\xi)Y\},
 \qquad
 \dot Y=-Y\{\xi B(\xi)+\zeta(\xi)Y\},
 \label{deg:tail-field}
\end{equation}
where \(\epsilon>0\) and \(B,Q\) are positive units, with transverse sections
\begin{equation}
 Y=\ell_0P(\xi),\qquad 0<\ell_0\ll1,
 \label{deg:moving-section}
\end{equation}
Put
\[
 A(\xi)=\int^\xi {xB(x)\over P(x)}\,dx,
 \qquad \mathcal D={P(\xi)\over \xi B(\xi)}\partial_\xi .
\]

\begin{lemma}[Exact moving-section identities]\label{deg:lem-moving-lip}
On a positive graph of \eqref{deg:tail-field} with
\(P>0\) and \(P+QY>0\), put \(W=Y/P\). Then
\begin{equation}
 {d\over d\xi}\left\{
  \epsilon\log Y-A(\xi)\right\}
 ={W\over1+QW}\left\{\zeta-Q{\xi B\over P}\right\}.
 \label{deg:exact-action-cocycle}
\end{equation}
If this graph joins the two lips at \(u\) and \(-r\), normalize
\(A(0)=0\) and write \(A_+(u)=A(u)\), \(A_-(r)=A(-r)\).
With \(P_+(u)=P(u)\), \(P_-(r)=P(-r)\), one has
\begin{equation}
 A_+(u)-A_-(r)
 =\epsilon\{\log P_+(u)-\log P_-(r)+R(u,r)\}.
 \label{deg:action-matching}
\end{equation}
where \(\epsilon R\) is the integral of the right-hand side of
\eqref{deg:exact-action-cocycle} from \(u\) to \(-r\) along this
graph. The two endpoint polynomials and the whole integral are retained.

For \(f=dY/d\xi\), its nonlinear derivatives and the inward
action density \(dq=|dA|/\epsilon\), away from \(\xi=0\), are
\begin{equation}
 \partial_Y^nf=(-1)^{n-2}n!
 \frac{PQ^{n-2}(\zeta P-\xi BQ)}{\epsilon(P+QY)^{n+1}},
 \quad n\ge2,\qquad
 |d\xi|=\frac{\epsilon P}{|\xi|B}\,dq.
 \label{deg:bell-density}
\end{equation}
For a moving initial curve \((\xi_l(x),Y_l(x))\), let
\(Y_T(x)\) be its graph value at a fixed terminal abscissa and let
\(\varphi_T\) be a fixed regular terminal adapter. Write
\(J\) for the scalar first variation. Then
\begin{equation}
 \frac d{dx}\varphi_T(Y_T(x))
 =\varphi_T'(Y_T)J(\xi_T,\xi_l)
                 \{Y_l'(x)-f(\xi_l,Y_l)\xi_l'(x)\}.
 \label{deg:moving-first-variation}
\end{equation}
In the label \(b=-x^2/2\), this derivative is multiplied by
\(-1/x\). Coefficient-only affine calibrations leave logarithmic
slope derivatives unchanged. The quantitative bounds for the actual
polynomial passages are established separately below and in
Appendix~\ref{app:moving-section-bookkeeping}.
\end{lemma}

\begin{proof}
Division of the graph equation and integration give the first two
identities; rational differentiation gives \eqref{deg:bell-density}.
The initial variation at fixed abscissa is
\(Y_l'-f_l\xi_l'\); its scalar variational equation gives
\eqref{deg:moving-first-variation}. The remaining statements are
the ordinary chain rule.
\end{proof}

\subsection{A compact-spectrum local Dulac estimate}

The following local Dulac estimate is used in the compact positive-ratio
passage and in the stopped determinant chart. Its logarithmic remainder
bounds remain uniform when the characteristic ratio crosses one.

\begin{lemma}[Uniform local positive Dulac estimate]
\label{lem:uniform-positive-dulac}
Consider a compact analytic family of hyperbolic saddles with fixed compatible
entrance and exit sections.  Suppose the positive tangential and normal rates
\(L,\nu\) are bounded above and away from zero and
\mbox{\(\rho=\nu/L\in[\rho_0,\rho_1]\Subset(0,\infty)\)}. Writing \(\theta\) for
the family parameter, in a positive local coordinate \(H\) which vanishes at
the incoming separatrix, the local Dulac map has the form
\[
 \mathcal D(H)=c_\theta H^\rho e^{\mathcal R_\theta(H)},
 \qquad c_\theta>0.
\]
For every \(0<\epsilon<\min(1,\rho_0)\), there is a uniform constant
\(C<\infty\) such that
\[
 \left|(H\partial_H)^j\mathcal R_\theta(H)\right|
 \le C H^\epsilon(1+|\log H|),
 \qquad 0\le j\le2.
\]
The estimate includes \(\rho=1\), and all local coordinate changes preserve
the positive-time orientation.
\end{lemma}
\begin{proof}
Straighten the invariant axes. A positive time factor gives
\(u'=Lu\). Linearize the one-dimensional stable-axis equation and
then replace \(v\) by \(ve^{h(u)}\), where
\(h'(u)=(b(u,0)-\nu)/(Lu)\). The quotient is removable at zero.
The remaining normal-coefficient defect vanishes on both axes, so
on one common small rectangle the equation is
\[
 u'=Lu,\qquad v'=-\nu v(1+uvG_\theta(u,v)),
 \qquad \rho=\nu/L\in[\rho_0,\rho_1].
\]
These changes use no division by \(\rho-1\), preserve positive
time, and have common finite derivative bounds on the compact family.
Take the entrance \((H,v_0)\) and the exit \(u=u_0\).
All phase derivatives below hold the coefficients and these cuts fixed.
Put \(\mathsf B=-\rho vG_\theta\). The graph equation is
\[
 \partial_u\log v=-\rho/u+\mathsf B(u,v).
\]
After shrinking the rectangle, \(v\) decreases. Its integral
equation first gives \(v(u;H)\le C(H/u)^\rho\), and then
\[
 I_\rho(H):=H^\rho\int_H^{u_0}u^{-\rho}\,du
 \le CH^{\min(1,\rho)}(1+|\log H|).
\]
The bound is uniform through \(\rho=1\); no separate bound for
\((1-\rho)^{-1}\) is taken.

Let \(\mathcal E=H\partial_H\), \(P=\mathcal E\log v\), and
\(Z=\mathcal E^2\log v\), at fixed \(u\). Differentiating the
moving initial condition \(v(H;H)=v_0\), as in
\eqref{deg:scalar-moving-bell}, gives
\[
 \begin{aligned}
 P(H)&=\rho-H\mathsf B(H,v_0),& Z(H)&=O(H),\\
 P_u&=\mathsf B_{\log v}P,
 &Z_u&=\mathsf B_{\log v}Z+\mathsf B_{\log v\log v}P^2.
 \end{aligned}
\]
The two logarithmic height derivatives of \(\mathsf B=-\rho vG\)
are \(O(v)\). Variation of constants therefore bounds \(P>0\) and gives
\( |P(u_0)-\rho|+|Z(u_0)|\le C(H+I_\rho(H))\).
For
\[
 \mathcal R(H)=\log\frac{v(u_0;H)}{v_0(H/u_0)^\rho},
\]
one has \(\mathcal E\mathcal R=P(u_0)-\rho\) and
\(\mathcal E^2\mathcal R=Z(u_0)\); its value is bounded by
the integral of \(\mathsf B\). These prove the stated estimate
for every fixed \(0<\epsilon<\min(1,\rho_0)\).
The endpoint \(H=0\) is a separatrix endpoint, not a through orbit.
\end{proof}
